\documentclass[11pt,a4paper]{letter}
\usepackage[utf8]{inputenc}
\usepackage[T1]{fontenc}
\usepackage[french]{babel}
\usepackage[left=2.5cm,right=2.5cm,top=2cm,bottom=2cm]{geometry}
\usepackage{hyperref}
\usepackage{xcolor}
\usepackage{fontawesome5}

\definecolor{primary}{RGB}{35, 55, 123}
\hypersetup{colorlinks=true, urlcolor=primary}

\pagestyle{empty}

\begin{document}

% --- EN-TÊTE EXPÉDITEUR ---
\begin{flushleft}
\textbf{Loïc Aron MBASSI EWOLO} \\
Élève-ingénieur ENSEA (2A) \\
Cergy, France \\
\faPhone\ (+33) 7 59 52 33 98 \\
\faEnvelope\ \href{mailto:loicaron.mbassiewolo@ensea.fr}{loicaron.mbassiewolo@ensea.fr} \\
\faLinkedin\ \href{https://www.linkedin.com/in/loic-mbassi}{linkedin.com/in/loic-mbassi}
\end{flushleft}

\vspace{0.3cm}

% --- DESTINATAIRE ---
\begin{flushright}
\textbf{ValoTec} \\
Service Recrutement / R\&D \\
Villejuif Bio Park \\
1 mail du Professeur Georges Mathé \\
94800 Villejuif \\
\today
\end{flushright}

\vspace{0.3cm}

% --- OBJET ---
\noindent\textbf{Objet :} Candidature Stage Ingénieur R\&D Électronique (Avril 2026) + Projet Alternance

\vspace{0.4cm}

% --- CORPS DE LA LETTRE ---
Madame, Monsieur,

\vspace{0.3cm}

L'intersection entre l'ingénierie électronique et le domaine médical est le fil conducteur que j'ai choisi pour ma carrière. Votre positionnement au sein du Villejuif Bio Park et votre expertise dans le développement de dispositifs médicaux actifs font de ValoTec l'environnement idéal pour concrétiser cette ambition. Actuellement en double diplôme ingénieur à l'ENSEA (Cergy) et à l'École Polytechnique de Yaoundé, je souhaite mettre ma rigueur scientifique et mon savoir-faire technique au service de vos projets R\&D.

\vspace{0.3cm}

Mon profil se distingue par une approche concrète de l'ingénierie électronique :

\begin{itemize}
    \item \textbf{Habileté manuelle :} Au laboratoire IOTA, j'ai contribué au projet Harmony Gloves (gant connecté pour la langue des signes). J'ai réalisé la \textbf{soudure des composants} (capteurs flex, IMU) et participé à l'intégration mécanique du prototype.
    \item \textbf{Instrumentation :} Validation des signaux à l'\textbf{oscilloscope}, calibration des capteurs, tests fonctionnels sur banc.
    \item \textbf{Programmation embarquée :} Développement en \textbf{C} sur STM32, maîtrise du \textbf{VHDL} pour conception FPGA, protocoles I2C/SPI.
    \item \textbf{CAO :} Familiarité avec \textbf{Altium Designer} et KiCad pour le routage PCB.
\end{itemize}

\vspace{0.2cm}

En parallèle, je cultive un bagage algorithmique solide. Participant au challenge robotique CoHoMa (Défense), j'y implémente des algorithmes de SLAM sur cibles Nvidia Jetson. Mon passage au Mila (Institut Québécois d'IA) m'a apporté la rigueur scientifique nécessaire à la rédaction de rapports techniques et à l'expérimentation méthodique.

\vspace{0.3cm}

Ce qui m'attire chez ValoTec : la pluridisciplinarité des projets, le travail en binôme/trinôme avec revue systématique, et la proximité avec des clients Medtech exigeants. Votre approche de formation par l'expérimentation correspond exactement à ma manière d'apprendre.

\vspace{0.3cm}

Souhaitant m'établir et poursuivre ma carrière d'ingénieur en France, je suis disponible dès \textbf{avril 2026} pour un stage de 4 à 5 mois. Mon objectif est de m'investir pleinement dans vos projets pour prouver ma valeur et poursuivre notre collaboration via un \textbf{Contrat de Professionnalisation} pour ma dernière année à l'ENSEA (dès septembre 2026).

\vspace{0.3cm}

Je serais ravi de vous présenter mes prototypes et de discuter de vos enjeux d'industrialisation lors d'un entretien.

\vspace{0.4cm}

Je vous prie d'agréer, Madame, Monsieur, l'expression de mes salutations distinguées.

\vspace{0.8cm}

\noindent\textbf{Loïc Aron MBASSI EWOLO} \\
\textit{Élève-Ingénieur ENSEA / Polytechnique Yaoundé}

\end{document}
